\section{Electron cascades}\label{app:E.24}

For most of the materials ranked in \cref{chap:17} no first-principles cascade exists. The method is the cascade of \cref{app:E.13}; \cref{tab:E24-summary} collects the grids, tests and costs of the calculations that apply it.

\subsection{Silicon, germanium and GaAs}

The output is the distribution $P_n(\vect{q}, \omega, T) = P(N_{\mathrm{eh}} = n|\vect{q}, \omega, T)$ of the pair number for a deposit of energy $\omega$ at momentum transfer $\vect{q}$, the distribution that \cref{sec:18.1} builds from the branching process and \cref{sec:18.4} compares with measurements. Its moments are $\varepsilon_{\mathrm{eh}}(\vect{q}, \omega, T)$ and $F_{\mathrm{eh}}(\vect{q}, \omega, T)$ of \cref{eq:E.13-moments}, and the variance is also reported split into the terms $\mathcal{V}_<$, $\mathcal{V}_>$ and $\mathcal{V}_{<>}$ of \cref{sec:18.2}, from below and above the impact-ionization threshold and their covariance. For silicon the run also records the energies of the carriers at the end of the cascade and the separations of the pairs they form, which \cref{sec:18.5} compares with the field-separation radius $r_{\mathcal{E}}$ and over which the capture probability of the last subsection is averaged. The grid is $\omega$ from $\Eg$ to \SI{50}{eV} in steps of \SI{50}{meV}, $|\vect{q}|$ from 0 to $25\alpha m_e$ in steps of $\alpha m_e$ along three directions plus the angular average, and five temperatures, 0, 50, 100, 150 and \SI{300}{K}. Above \SI{50}{eV} the moments are matched to their asymptotic forms, a mean $\omega/\varepsilon_{\mathrm{eh}}$ plus a constant offset and a variance $F_{\mathrm{eh}}\omega/\varepsilon_{\mathrm{eh}}$. On $\mathfrak{T}$ (\cref{sec:8.5}) the bands are $G_0W_0$ quasiparticles, the screening is RPA with local fields ($\mathcal{X}$ = RPA), the phonon kernel is at $n_g = 2$ and $n_B = 0$, and the events are incoherent (\cref{app:E.13}). Temperature enters three times: through the occupation factors of the phonon kernel, through the Bose factor $1 + n_{\mathrm{B}}(\omega)$ of the pair or plasmon excited at energy transfer $\omega$ in the impact-ionization kernel, and through the temperature-renormalized gap (\cref{app:E.8}).

The bands come from Quantum ESPRESSO and $G_0W_0$ in BerkeleyGW~\cite{deslippe2012berkeleygw} or Yambo~\cite{sangalli2019many}, Wannier-interpolated up to about \SI{15}{eV} above the conduction minimum. Silicon needs 16 to 20 disentangled $sp^3d$ Wannier functions; germanium and GaAs need about as many, with spin--orbit coupling. A germanium band structure enters only if it passes the check of \cref{app:E.14}: L below $\mathrm{\Gamma}$, both gaps within \SI{0.1}{eV} of the measured ones, and the transverse mass at L and the mass at $\mathrm{\Gamma}$ against measurement. The germanium bands are $G_0W_0$ with spin--orbit coupling on a PBE or hybrid ground state. Above \SI{15}{eV} the states are taken as free-electron-like with $GW$ lifetimes, the range in which the free-particle form of the single-parameter model is accurate, and the matching point serves as a check. The screened interaction is the full-frequency $\varepsilon^{-1}_{\vect{G}\vect{G}'}(\vect{p}, \omega)$ for $\omega$ from 0 to \SI{60}{eV}, on about 200 frequencies that are denser below \SI{20}{eV}; carriers above \SI{60}{eV}, up to the energy cap, use the free-electron kernel of \cref{app:E.13}, so the grid of $W^{\mathrm{scr}}$ need not reach the cap. The screened interaction uses $10^3$ to $12^3$ momenta $\vect{p}$, $N_G \approx 200$ to 300 and at least 200 empty bands, and $\Imag\varepsilon^{-1}$ at the plasmon, whose loss peak lies at \SI{16.4}{eV} in the optical loss function of silicon (\cref{sec:4.4}), is converged to \SI{2}{\percent}. For GaAs the total dielectric function includes the ionic part built from the measured $\varepsilon_\infty$ and the measured LO--TO splitting, which gives $Z^*/\varepsilon_\infty = 0.194$ (\cref{sec:5.5,sec:11.9}); a semilocal $\varepsilon_\infty$ of GaAs is not a safe input. The impact-ionization kernel of \cref{eq:E.13-imsigma} is tabulated with local fields on iso-energy surfaces, about $10^3$ $\vect{k}$ points per \SI{50}{meV} bin. The electron--phonon kernel comes from EPW~\cite{ponce2016epw}, with $8^3$ coarse grids, fine $\vect{k}$ grids of $80^3$ to $100^3$ and $50^3$ phonon wavevectors sampled by Monte Carlo. GaAs needs the long-range dipole and quadrupole terms, and every material needs the same vertex and phonons as the response code (\cref{app:E.5}).

The pair number is computed in two stages. The first generation, the primary carriers and the products of their first impact ionization, keeps its full $\vect{k}$ resolution, because the reduction to energy alone is least accurate there; it is seeded with the final states of the dark-matter response (\cref{sec:18.3}). Later generations use the energy-only reduction of \cref{sec:18.1}, solved deterministically by the moment equations on an energy grid of bin $h$, with kernels averaged over iso-energy surfaces. For $N_s$ sampled first steps with subtree means $M_i$ and variances $V_i$, the estimator is $\bar n = 1 + \operatorname{mean}_iM_i$ and $\operatorname{Var}N = \operatorname{mean}_iV_i + \operatorname{var}_iM_i$, the law of total variance, with errors from 20 batch means. The relative error of $\operatorname{var}_iM_i$ is $(2/N_s)^{1/2}$ for a near-Gaussian $M$, so a Fano factor to \SI{0.5}{\percent} needs at most $2/0.005^2 = 8\times10^4$ samples; the run starts at $10^4$ and doubles until the batch errors fall below \SI{0.5}{\percent} on $F_{\mathrm{eh}}$ and \SI{0.1}{\percent} on $\varepsilon_{\mathrm{eh}}$. The plasmon part of $\Imag W^{\mathrm{scr}}$ is tagged and its share of the variance reported (\cref{app:E.13}).

The grid is converged when the bin, the energy cap, the matching point, the fit window and the $\vect{k}$ mesh each pass. Halving $h$ from at most \SI{10}{meV} changes $\varepsilon_{\mathrm{eh}}$ by less than $10^{-3}$ and $F_{\mathrm{eh}}$ by less than $3\times10^{-3}$, relative. Raising the energy cap from 100 to \SI{150}{eV}, and the free-electron matching point from 15 to \SI{20}{eV}, change both by less than $10^{-3}$. The asymptotic slopes fitted over $[X/2, X]$ and over $[3X/4, X]$, for a top energy $X$, agree to $10^{-4}$. A second $\vect{k}$ mesh for the iso-energy averages changes $\varepsilon_{\mathrm{eh}}$ and $F_{\mathrm{eh}}$ by less than $10^{-3}$. Energy is conserved along every path to $10^{-12}$, and $F_{\mathrm{eh}}$ computed from the pair number and from the heat agree to $10^{-9}$. Before the full kernels are switched on, the code must reproduce, with free-particle kernels in their place, the single-parameter model of \cref{app:E.13} within its Monte Carlo scatter of about \SI{2}{\percent}. The validation, not used in any fit, is the silicon measurement $\varepsilon_{\mathrm{eh}} = \SI{3.749 \pm 0.001}{eV}$ and $F_{\mathrm{eh}} = 0.119 \pm 0.002$ at \SI{123}{K}~\cite{rodrigues2021absolute} and the silicon quantum yield of photons between 3 and \SI{8}{eV} and between 50 and \SI{1500}{eV}~\cite{ramanathan2020ionization}, with the 50 to \SI{1500}{eV} data run at \SI{300}{K}, their measurement temperature. For germanium no measurement at \SI{77}{K} uses the same definition of the pair energy; until one does, the germanium value is an input, not a test. A photon seed at $q \approx 0$ and a dark-matter seed at $q$ of 1 to \SI{10}{keV}, run in one code, give the $q$ dependence of the yield directly.

The cost per material is $10^4$ to $10^5$ core-hours for the screened interaction, as much for the impact-ionization kernel, and about $10^4$ for the phonon kernel; the three materials together need $10^5$ to $6\times10^5$ core-hours. Memory is set by $W^{\mathrm{scr}}$: at $N_G = 300$, 200 frequencies and the 72 irreducible points of a $12^3$ grid, $72\times300^2\times200\times\SI{16}{byte} \approx \SI{21}{GB}$. Beyond the pitfalls of \cref{app:E.13}, each phonon step must be an integer number of energy bins to better than \SI{1}{\percent}, either by a bin commensurate with the dominant phonon or by interpolating the other branches, since an off-grid step biases every phonon loss the same way. The quasiparticle residue enters once, in the rate or in the weight, never in both. Kohn--Sham and quasiparticle energies must not be mixed between kernels. $F_{\mathrm{eh}}$ is a small difference of large terms (\cref{sec:18.2}), so the three variance terms are reported with their errors. The results replace the model values of \cref{sec:18.4} for the three materials and fix the $q$ dependence of the yield used in \cref{sec:18.1}.

\subsection{Wide-gap and polar crystals}

Sapphire, CaWO$_4$, LiF and NaF, among the materials ranked in \cref{chap:17}, have no measured Fano factor and no validated pair energy, and \cref{sec:18.4} shows that they lie outside the range in which the single-parameter model was validated. These four come first; 4H-SiC, diamond, 3C-SiC, GaN, AlN and TeO$_2$ follow. The quantities and grids are those of the silicon cascade. Polar crystals add the long-range Fr\"ohlich vertex, built with the same Born charges and dielectric constants as the response code (\cref{app:E.4,app:E.5}). LiF and NaF have excitons bound by an eV or more, so the Bethe--Salpeter equation of \cref{app:E.6} is solved as well. The bookkeeping must then state whether a bound exciton counts as a pair. A carrier between $\Eg - E_x^b$ and the ionization threshold can also create a bound pair directly, and the cascade must include that channel. Tungsten in CaWO$_4$ needs spin--orbit coupling.

The cells are small, two atoms for LiF and NaF, 10 for sapphire and 12 for the primitive cells of CaWO$_4$ and TeO$_2$, but the oxides have many bands. The cost is estimated at about $10^5$ core-hours for each oxide, $10^4$ to $10^5$ for each halide plus the Bethe--Salpeter step, and twice the silicon cost for diamond, the silicon carbides, GaN and AlN; the set needs $6\times10^5$ to $2.5\times10^6$ core-hours. The silicon validation comes first, because silicon alone has a measurement at the precision that tests a Fano factor, and the other materials are run after the silicon validation passes, reusing the silicon kernels. Their benchmark is the full-band Monte Carlo with density-functional inputs, which gives \SI{7.3}{eV} against the measured \SI{7.8}{eV} for 4H-SiC, \SI{6.6}{eV} against \SI{6.9}{eV} for 3C-SiC and \SI{16.7}{eV} against \SI{18.0}{eV} for SiO$_2$~\cite[Table~II]{nielsen2026first}, a few to \SI{15}{\percent} below measurement where measurements exist. The computed pair energies and the first ab initio Fano factors of these materials replace the estimates of \cref{sec:18.4}.

\subsection{Impact ionization beyond first order}

\Cref{app:E.13} gives two definitions of the impact-ionization rate, from the imaginary part of the $GW$ self-energy with its residue $z_{\mathrm{qp}}$ and from the golden rule with exchange, which differ by a factor between 0.5 and 1. A constant factor does not matter once a model is calibrated on silicon. In the single-parameter model a factor $z_{\mathrm{qp}} = 0.8$ on the rate raises the fitted phonon-to-ionization constant $A_{\mathrm{AB}} = \SI{5.703}{eV^3}$ by \SI{25}{\percent}, $\Delta A_{\mathrm{AB}} = \SI{1.43}{eV^3}$, and with the secant slope $\partial\varepsilon_{\mathrm{eh}}/\partial A_{\mathrm{AB}} = \SI{0.076}{eV^{-2}}$ of \cref{sec:18.4} it raises $\varepsilon_{\mathrm{eh}}$ by \SI{0.11}{eV}, \SI{2.9}{\percent}, which recalibration removes. A factor that varies with energy cannot be removed, and it changes $F_{\mathrm{eh}}$, which depends on how the rate rises above threshold (\cref{sec:18.5}). We compute the ratio $r_{\mathrm{ii}}(E)$ of the impact-ionization rate beyond first order in $W^{\mathrm{scr}}$ to the first-order rate, on \SI{50}{meV} bins from threshold to about \SI{10}{eV}, for silicon.

The calculation evaluates $\Imag\Sigma_{n\vect{k}}$ in $G_0W_0$ with and without $z_{\mathrm{qp}}$, and adds the second-order screened exchange as the first vertex correction beyond $GW$. Vertex corrections of this kind shift the absolute band energies of solids by tenths of an eV~\cite{grueneis2014ionization}, and the question here is how much they change the rate. The grids are those of the cascade: $N_G \approx 250$, 200 empty bands, the full-frequency $\varepsilon^{-1}$ and $10^3$ $\vect{k}$ points per bin. The ratio is converged when it changes by less than \SI{2}{\percent} per bin with $N_G$ raised to 350, the bands to 400 and the $\vect{k}$ points to $2\times10^3$; the second-order exchange is reported separately to the same precision. A variation below \SI{2}{\percent} across the window counts as constant. The cascade is then rerun with $\varGamma_{\mathrm{ii}} \to r_{\mathrm{ii}}(E)\varGamma_{\mathrm{ii}}$, and the change of $F_{\mathrm{eh}}$, for about $3\times10^4$ core-hours in all, is the uncertainty that the order of the rate adds to $F_{\mathrm{eh}}$ in \cref{sec:18.5}.

\subsection{Momentum memory between events}

The energy-only reduction of \cref{sec:18.1} assumes that a carrier forgets its momentum on its iso-energy surface between two events that change its energy. Suppose momentum-randomizing scattering, by intervalley and acoustic phonons, is a Poisson process of rate $\varGamma_k$ that leaves the carrier uniform on the surface, while energy-changing events occur at rate $\varGamma_E = \varGamma_{\mathrm{ii}} + \varGamma_{\mathrm{ph,loss}}$. The probability that the next event changes the energy before any randomization is
\begin{equation}
  w_{\mathrm{mem}}(E) = \frac{\varGamma_E}{\varGamma_E + \varGamma_k} ,
  \label{eq:E.24-wmem}
\end{equation}
the ratio of competing Poisson rates. At the next energy-changing event the momentum distribution is then $1 - w_{\mathrm{mem}}$ times the uniform one plus $w_{\mathrm{mem}}$ times the remembered one. For any cascade observable $m_{\vect{k}}(E)$ that depends on the starting $\vect{k}$, such as the mean pair number or its square, the error of the reduction per generation is therefore at most $w_{\mathrm{mem}}(E)$ times the spread of $m_{\vect{k}}$ over the surface, $\max_{\vect{k}}m_{\vect{k}} - \min_{\vect{k}}m_{\vect{k}}$. Within about \SI{1}{eV} of threshold, intervalley scattering at $10^{13}$ to \SI{e14}{s^{-1}} dominates and $w_{\mathrm{mem}}$ is small. Above 3 to \SI{5}{eV} in silicon, impact ionization outpaces the phonons and $w_{\mathrm{mem}}$ can reach one half, so the bound is useful only if the spread is small.

The calculation gives the spread of the mean pair number and of its second moment, and $w_{\mathrm{mem}}(E)$, for silicon between 3 and \SI{8}{eV} on \SI{50}{meV} bins. They come from the $\vect{k}$-resolved second generation of the cascade, with at least 200 $\vect{k}$ samples per bin, weighted by area on the surface, and at least $10^3$ cascades per sample. $\varGamma_k$ comes from EPW, $\varGamma_E$ from the impact-ionization and phonon rates of the same run. The spread is converged to \SI{10}{\percent} when the $\vect{k}$ samples are doubled to 400 and the cascades quadrupled, for $10^3$ to $10^4$ core-hours. The maximum minus the minimum of $n_k$ noisy estimates is biased upward by about $2\sigma(2\ln n_k)^{1/2}$, with $\sigma$ the noise of one estimate, so the range between the 5th and 95th percentiles is used, or the noise is subtracted. The bound rests on the assumption that randomization leaves the carrier uniform on the surface, itself a form of the energy-only reduction; it bounds the error under that assumption and is not an independent test of it.

\subsection{A library of pair-number distributions}

The response libraries of \cref{sec:21.6} need, for every final-state class $\Lambda$ of every material, the full distribution $P_{\mathrm{casc}}(N_{\mathrm{eh}}|\Lambda)$ (\cref{sec:18.3}). The reduced cascade gives it cheaply. With energy-only branching, random final momenta for impact ionization, one phonon energy $\omega_{\mathrm{ph}}$ and one phonon-to-ionization constant, the generating function $\mathcal{G}_\Lambda(z) = \sum_Nz^NP_{\mathrm{casc}}(N|\Lambda)$ obeys the recursion of \cref{sec:18.1} in the energies $x_e$ and $x_h$ of the seed electron and hole above their band edges. We evaluate $\mathcal{G}_\Lambda$ at the $N_z$ roots of unity $z_j = \ee^{2\pi\ii j/N_z}$. The discrete Fourier sum
\begin{equation}
  P_{\mathrm{casc}}(N|\Lambda) = \frac{1}{N_z}\sum_{j=0}^{N_z-1}z_j^{-N}\,\mathcal{G}_\Lambda(z_j) ,
  \label{eq:E.24-dft}
\end{equation}
returns $\sum_{N'}P_{\mathrm{casc}}(N + N'N_z|\Lambda)$. This sum equals the exact probability as long as $N_z$ exceeds every pair number with appreciable weight, so $N_z$ is taken as the smallest power of two above $\bar N + 8\sigma_N$, with $\sigma_N$ the standard deviation of the pair number.

The energy grid has a step $h$ of about $\Eg/50$, chosen so that it divides $\Eg$, $\theta_{\mathrm{ii}}\Eg$ and $\omega_{\mathrm{ph}}$. With phonons the discretization error falls as $h^{3/2}$, so the values at $h$ and $h/2$ are extrapolated to zero step as $f(h/2) + [f(h/2) - f(h)]/(2^{3/2} - 1)$ (\cref{app:D.43}); the $h^2$ formula would leave a bias. The energies run to $20\Eg$, and beyond the table the distribution is the asymptotic lattice Gaussian of mean $\bar N$ and variance $\sigma_N^2$, with the Edgeworth terms added. The Gaussian is sampled at the integers, since integrating it over unit bins adds $1/12$ to the variance. The Gaussian is approached slowly: with $x$ the carrier energy in units of $\Eg$, the absolute pointwise error falls as $C/x$ and the relative skewness as $x^{-1/2}$, with $C \approx 0.24$ for silicon, so the match at $20\Eg$ carries an error of about $1.2\times10^{-2}$ in each probability. The tests are internal: the probabilities sum to one to machine precision and are non-negative; the derivatives of $\mathcal{G}_\Lambda$ at $z = 1$ reproduce the moment recursion to $10^{-9}$; and below the first impact-ionization threshold $P_{\mathrm{casc}}(N|\Lambda) = \delta_{N1}$. The reader that interpolates the tables must return that exact value below $\theta_{\mathrm{ii}}\Eg$ and interpolate only from the threshold upward, with a table row at the threshold itself: for $\theta_{\mathrm{ii}} > 1$ the distribution jumps there, and an interpolation across the jump puts weight on extra pairs below threshold.

The same tables carry the marks $\mathfrak{m}$ of \cref{eq:18.3-obs} that a readout needs (\cref{sec:21.6}): the phonon spectrum and the separate electron and hole counts, which follow from the same recursion with one further generating variable for each; the carrier positions, taken from the ranges between events in the $\vect{k}$-resolved cascade of silicon; and the stored energy, taken from the table of configurational energies of \cref{app:E.23}. A table takes seconds per material and parameter set, so five parameter sets for every ranked material take minutes on one core. The phonon step must be commensurate with the energy grid. Rounding $\omega_{\mathrm{ph}}$ to the grid changes it by up to $h/2$; with $h = 0.02\Eg$ in silicon that is up to \SI{12}{meV} of a \SI{63}{meV} phonon, and every phonon loss of the cascade inherits the error. Two smaller errors are aliasing in $z$, which \cref{eq:E.24-dft} avoids only if $N_z$ is large enough, and the seed dependence of the constant offset of $\bar N$. A seed that gives all the energy to one carrier and a seed that shares it have different offsets. Temperature beyond its effect on the gap, and excitons, are outside the reduced model. When the full-band kernels of the silicon cascade exist, they replace the free-particle forms in the same recursion.

\subsection{Exciton and capture parameters}

At millikelvin the pairs at the end of the cascade bind into free excitons, bind to impurities, are trapped, or recombine, and \cref{sec:18.5} writes these processes as rate equations whose rates are unknown in silicon and germanium. The exciton itself needs its binding energy $E_x^b$, radius $a_x$ and oscillator strengths, and its binding to neutral donors and acceptors. Capture needs the cross-section $\sigma_{\mathrm{cap}}(E, T, \mathcal{E}_{\mathrm{b}})$ of a carrier of energy $E$ on shallow centers in a bias field $\mathcal{E}_{\mathrm{b}}$, and multiphonon capture on deep levels. The kinetics needs the rates of exciton formation, trapping of free carriers, binding of excitons to each impurity species and decay of bound excitons, the luminescence fractions, the bimolecular formation coefficient, exciton--exciton annihilation and Auger recombination, all at 10, 20 and \SI{50}{mK}. The targets are silicon and germanium, then LiF, NaF, sapphire and CaWO$_4$.

On $\mathfrak{T}$ (\cref{sec:8.5}) the exciton parameters are at $\mathcal{X}$ = BSE on quasiparticle bands, the bound excitons at effective-mass level with a density-functional central cell, and the capture rates at $n_g = 2$. The exciton parameters come from $G_0W_0$ and the Bethe--Salpeter equation (\cref{app:E.6}) in BerkeleyGW or Yambo, which in silicon and germanium (for germanium on quasiparticle bands that pass the band check) checks the effective-mass binding of \cref{app:E.15}, with 4 to 8 valence and 4 to 8 conduction bands. In silicon and germanium the exciton radius spans many cells, and the $\vect{k}$ grid must resolve it: a spacing below $1/(3a_x)$ means $10^4$ to $10^5$ points, reached by double-grid interpolation, against $8^3$ to $12^3$ for the ionic crystals. When $E_x^b$ approaches the longitudinal optical phonon energy, the lattice part of the screening enters through a polaronic correction. Donor- and acceptor-bound excitons come from effective-mass theory with central-cell corrections from density-functional calculations in 216- to 512-atom cells. Capture on shallow centers follows cascade-capture theory~\cite{lax1960cascade}, a carrier captured if it emits phonons inside the Coulomb radius, with acoustic and optical emission rates from EPW and the vertex of the cascade. Deep levels use charge-transition levels and Huang--Rhys factors (\cref{app:E.10,app:E.22})~\cite{alkauskas2014first}. The Bethe--Salpeter step costs $10^3$ to $10^4$ core-hours for the semiconductors and about $10^4$ for the oxides; capture about $10^5$ core-hours per material; the six materials about $6.5\times10^5$.

The exciton levels of the ionic crystals are converged to \SI{1}{meV} when the $\vect{k}$ mesh goes from $8^3$ to $12^3$, and those of silicon and germanium when the double-grid density near the band edges is doubled, with the bands from four to six on each side; central-cell binding energies to \SI{10}{\percent} between 216 and 512 atoms. A rate that changes by more than \SI{10}{\percent} between 10 and \SI{50}{mK} is reported as a temperature dependence. The measured control is the free-exciton binding of silicon, \SI{14.7}{meV}~\cite{shaklee1970valley}. Formation of an exciton from two free carriers moves gap energy between pairs and excitons and releases only $E_x^b$, while exciton--exciton annihilation and Auger recombination remove gap energy, so the bimolecular channel must be named. The excitons cool by acoustic phonons slowly at millikelvin, so their temperature is not the lattice temperature. Capture at the end of the cascade changes the pair statistics. A capture probability shared by every carrier changes $\varepsilon_{\mathrm{eh}}$ and $F_{\mathrm{eh}}$ by \cref{eq:18.5-thin}. A capture probability $f(x)$ that depends on the final energy $x$ of the carrier is instead averaged, with the law of total variance, over the end-of-cascade energy distribution that the silicon cascade records. The parameters close the rate equations and the end-of-cascade populations of \cref{sec:18.5}.

\begin{table}[t]
  \caption{Electron-cascade calculations, with estimated costs in core-hours. GW+BSE, quasiparticle and Bethe--Salpeter calculation.}
  \label{tab:E24-summary}
  \centering
  \small
  \begin{tabular}{p{0.19\linewidth}p{0.29\linewidth}p{0.27\linewidth}p{0.13\linewidth}}
    \toprule
    quantity & method and grid & acceptance & cost \\
    \midrule
    $P_n(\vect{q}, \omega, T)$, $\varepsilon_{\mathrm{eh}}$, $F_{\mathrm{eh}}$, end-of-cascade energies and pair separations (Si, Ge, GaAs) & $G_0W_0$ bands (Ge after the band check), full-frequency $W^{\mathrm{scr}}$ to \SI{60}{eV}, EPW; $k$-resolved first generation & Si \SI{3.749}{eV}, 0.119 at \SI{123}{K}; quantum yield at \SI{300}{K} & $10^5$--$6\times10^5$ \\
    $\varepsilon_{\mathrm{eh}}$, $F_{\mathrm{eh}}$ (wide-gap, polar) & as above, Fr\"ohlich vertex, BSE for halides & silicon validation first; DFT full-band benchmarks & $6\times10^5$--$2.5\times10^6$ \\
    $r_{\mathrm{ii}}(E)$, Si & $G_0W_0$ with and without $z_{\mathrm{qp}}$, plus second-order exchange & \SI{2}{\percent} per \SI{50}{meV} bin & $3\times10^4$ \\
    $w_{\mathrm{mem}}(E)$ and spread, Si, 3--\SI{8}{eV} & $k$-resolved second generation, 200 $\vect{k}$ per bin & \SI{10}{\percent}; percentile range & $10^3$--$10^4$ \\
    $P_{\mathrm{casc}}(N|\Lambda)$ library, marks & reduced recursion, $h = \Eg/50$, \cref{eq:E.24-dft}, $h^{3/2}$ extrapolation & sum to one; moments; $\delta_{N1}$ & $\le 1$ \\
    Exciton and capture set & GW+BSE, central-cell DFT, cascade capture & Si $E_x^b = \SI{14.7}{meV}$ & $6.5\times10^5$ \\
    \bottomrule
  \end{tabular}
\end{table}
